\documentclass[12pt,a4paper]{article}
\usepackage[utf8]{vietnam}
\usepackage{amsmath,amssymb}
\usepackage{tabularx}
\usepackage{tikz}
\usetikzlibrary{calc,angles,quotes,arrows.meta,patterns}
\usepackage{xcolor}
\usepackage{enumitem}
\usepackage{fancyhdr}
\usepackage[top=1.2cm, bottom=1.2cm, left=1.5cm, right=1.5cm]{geometry}
\usepackage[loigiai]{ex_test}
\SAOPTN{widthSA=2.6,heightSA=0.62}

\setlength{\headheight}{15pt}
\setlength{\parskip}{0pt}

\pagestyle{fancy}
\fancyhf{}
\rhead{\textit{Môn Toán 6 -- Hướng dẫn giải Đề số 10}}
\lhead{\textit{Bộ đề ôn tập Giữa học kì I}}
\cfoot{\thepage}

\begin{document}

\begin{center}
  \textbf{\Large HƯỚNG DẪN GIẢI CHI TIẾT VÀ BẢNG ĐÁP ÁN}\\
  \textbf{\large ĐỀ ÔN TẬP GIỮA HỌC KỲ I -- ĐỀ SỐ 10}\\
  \textit{Môn: Toán 6 -- Năm học 2026 -- 2027}
\end{center}
\vspace{2mm}

% --- BẢNG ĐÁP ÁN TRẮC NGHIỆM TÓM TẮT ---
\noindent\textbf{\large A. BẢNG ĐÁP ÁN TRẮC NGHIỆM TỔNG HỢP}
\vspace{1.5mm}

\noindent\textbf{1. Phần I (Mỗi câu trả lời đúng được 0,25 điểm):}
\begin{center}
\begin{tabular}{|c|c|c|c|c|c|c|c|c|c|c|c|c|}
\hline
\textbf{Câu} & \textbf{1} & \textbf{2} & \textbf{3} & \textbf{4} & \textbf{5} & \textbf{6} & \textbf{7} & \textbf{8} & \textbf{9} & \textbf{10} & \textbf{11} & \textbf{12} \\
\hline
\textbf{Đ/A} & \textbf{C} & \textbf{A} & \textbf{C} & \textbf{C} & \textbf{A} & \textbf{B} & \textbf{A} & \textbf{B} & \textbf{B} & \textbf{C} & \textbf{A} & \textbf{A} \\
\hline
\end{tabular}
\end{center}
\vspace{1mm}

\noindent\textbf{2. Phần II (Đúng 1 ý: 0,1đ; 2 ý: 0,25đ; 3 ý: 0,5đ; 4 ý: 1,0đ):}
\begin{center}
\begin{tabular}{|c|c|c|c|c||c|c|c|c|c|}
\hline
\textbf{Câu} & \textbf{a)} & \textbf{b)} & \textbf{c)} & \textbf{d)} & \textbf{Câu} & \textbf{a)} & \textbf{b)} & \textbf{c)} & \textbf{d)} \\
\hline
\textbf{Câu 1} & \textbf{Đ} & \textbf{S} & \textbf{Đ} & \textbf{Đ} & \textbf{Câu 2} & \textbf{Đ} & \textbf{Đ} & \textbf{S} & \textbf{Đ} \\
\hline
\end{tabular}
\end{center}
\vspace{1mm}

\noindent\textbf{3. Phần III (Mỗi câu trả lời đúng được 1,0 điểm):}
\begin{center}
\begin{tabular}{|c|c|c|}
\hline
\textbf{Câu} & \textbf{Câu 1} & \textbf{Câu 2} \\
\hline
\textbf{Đáp án} & \textbf{12} & \textbf{32} \\
\hline
\end{tabular}
\end{center}
\vspace{3mm}

% --- B. LỜI GIẢI CHI TIẾT ---
\noindent\textbf{\large B. LỜI GIẢI CHI TIẾT CÁC CÂU HỎI}
\vspace{2mm}

\noindent\textbf{\large PHẦN I. CÂU TRẮC NGHIỆM NHIỀU PHƯƠNG ÁN LỰA CHỌN}
\vspace{1mm}
\begin{ex}
Cho hai tập hợp $A = \{1; 2; 3; 4; 5\}$ và $B = \{a; b; c; d\}$. Có bao nhiêu tập hợp gồm hai phần tử thỏa mãn có một phần tử thuộc tập hợp $A$ và một phần tử thuộc tập hợp $B$?
\choice
{$16$}
{$18$}
{\True $20$}
{$9$}
\loigiai{
Tập $A$ có 5 phần tử, tập $B$ có 4 phần tử. Số cách chọn 1 phần tử thuộc $A$ là 5 cách, số cách chọn 1 phần tử thuộc $B$ là 4 cách. Số tập hợp thỏa mãn là: $5 \times 4 = 20$ (tập hợp).
}
\end{ex}

\begin{ex}
Tìm số tự nhiên $x$, biết $x + 20 = 42$:
\choice
{\True $22$}
{$62$}
{$2$}
{$18$}
\loigiai{
Ta có: $x = 42 - 20 = 22$.
}
\end{ex}

\begin{ex}
Số $308$ đọc là:
\choice
{Ba trăm tám mươi}
{Ba mươi tám}
{\True Ba trăm linh tám}
{Ba linh tám}
\loigiai{
Số $308$ có chữ số hàng trăm là 3, chữ số hàng chục là 0, chữ số hàng đơn vị là 8 nên được đọc là "Ba trăm linh tám".
}
\end{ex}

\begin{ex}
Hợp số là:
\choice
{Số tự nhiên có nhiều hơn 2 ước}
{Số tự nhiên có 4 ước}
{\True Số tự nhiên lớn hơn 1, có nhiều hơn 2 ước}
{Số tự nhiên lớn hơn 1, chỉ có hai ước là 1 và chính nó}
\loigiai{
Theo định nghĩa SGK, hợp số là số tự nhiên lớn hơn 1, có nhiều hơn 2 ước.
}
\end{ex}

\begin{ex}
Trong tam giác đều, mỗi góc có số đo bằng:
\choice
{\True $60^\circ$}
{$45^\circ$}
{$90^\circ$}
{$30^\circ$}
\loigiai{
Tam giác đều có 3 góc ở các đỉnh bằng nhau và mỗi góc bằng $60^\circ$.
}
\end{ex}

\begin{ex}
Khẳng định nào dưới đây là đúng?
\choice
{Hình bình hành có 4 góc bằng nhau}
{\True Hình lục giác đều có 3 đường chéo chính cùng đi qua một điểm}
{Tam giác đều là tam giác có 3 góc vuông}
{Hình chữ nhật có hai đường chéo bằng nhau và vuông góc với nhau}
\loigiai{
Hình lục giác đều có 3 đường chéo chính bằng nhau và cùng đi qua một điểm (tâm của lục giác đều).
}
\end{ex}

\begin{ex}
Kết quả của phép tính $7^{12} \cdot 7^2 \cdot 7^0$ dưới dạng một lũy thừa là:
\choice
{\True $7^{14}$}
{$7^{24}$}
{$7^{32}$}
{$7^7$}
\loigiai{
$7^{12} \cdot 7^2 \cdot 7^0 = 7^{12 + 2 + 0} = 7^{14}$.
}
\end{ex}

\begin{ex}
Cho hình bình hành $ABCD$ có $AB = 3\text{ cm}, BC = 2\text{ cm}$. Chu vi hình bình hành $ABCD$ là:
\choice
{$6\text{ cm}$}
{\True $10\text{ cm}$}
{$12\text{ cm}$}
{$5\text{ cm}$}
\loigiai{
Chu vi hình bình hành $ABCD$ là: $(3 + 2) \times 2 = 10\text{ (cm)}$.
}
\end{ex}

\begin{ex}
Tập hợp nào sau đây là tập hợp các số nguyên $\mathbb{Z}$?
\choice
{$\{1; 2; 3; 4; 5\}$}
{\True $\{\dots; -3; -2; -1; 0; 1; 2; 3; \dots\}$}
{$\{0; 1; 2; 3; 4; 5\}$}
{$\{0; -1; -2; -3; \dots\}$}
\loigiai{
Tập hợp các số nguyên gồm các số nguyên âm, số 0 và các số nguyên dương: $\mathbb{Z} = \{\dots; -3; -2; -1; 0; 1; 2; 3; \dots\}$.
}
\end{ex}

\begin{ex}
Số 18 được viết bằng chữ số La Mã là:
\choice
{XII}
{XV}
{\True XVIII}
{XX}
\loigiai{
$18 = 10 + 5 + 1 + 1 + 1$ tương ứng viết thành XVIII.
}
\end{ex}

\begin{ex}
Số nào sau đây là bội của $17$?
\choice
{\True $51$}
{$54$}
{$56$}
{$63$}
\loigiai{
Ta có $51 : 17 = 3$ (chia hết) nên 51 là bội của 17.
}
\end{ex}

\begin{ex}
Cho hình thoi có độ dài hai đường chéo là $2\text{ cm}$ và $4\text{ cm}$. Diện tích của hình thoi đó là:
\choice
{\True $4\text{ cm}^2$}
{$6\text{ cm}^2$}
{$8\text{ cm}^2$}
{$2\text{ cm}^2$}
\loigiai{
Diện tích hình thoi là: $S = \dfrac{2 \times 4}{2} = 4\text{ (cm}^2\text{)}$.
}
\end{ex}


\setcounter{ex}{0}
\vspace{2mm}
\noindent\textbf{\large PHẦN II. CÂU TRẮC NGHIỆM ĐÚNG SAI}
\vspace{1mm}
\begin{ex}
Cho tập hợp $B = \{x \in \mathbb{N} \mid 12 \le x < 18\}$. Xét tính đúng/sai của các khẳng định sau:
\choiceTF
{\True Tập hợp $B$ gồm các phần tử là $12; 13; 14; 15; 16; 17$}
{Số $18$ thuộc tập hợp $B$ (kí hiệu là $18 \in B$)}
{\True Trong tập hợp $B$, chỉ có duy nhất một số chia hết cho 5 là số $15$}
{\True Tổng của phần tử bé nhất và phần tử lớn nhất trong tập hợp $B$ là $29$}
\loigiai{
\textbf{a)} Đúng. Các số tự nhiên lớn hơn hoặc bằng 12 và nhỏ hơn 18 là 12, 13, 14, 15, 16, 17.

\textbf{b)} Sai. Vì điều kiện là $x < 18$ nên $18 \notin B$.

\textbf{c)} Đúng. Trong các phần tử của $B$, chỉ có số 15 có tận cùng là 5 nên chia hết cho 5.

\textbf{d)} Đúng. Phần tử bé nhất là 12, phần tử lớn nhất là 17. Tổng: $12 + 17 = 29$.
}
\end{ex}

\begin{ex}
Cho hình thang cân $ABCD$ có độ dài đáy nhỏ $AB = 4\text{ cm}$, đáy lớn $CD = 10\text{ cm}$ và chiều cao $h = 5\text{ cm}$. Xét tính đúng/sai của các khẳng định sau:
\choiceTF
{\True Hai cạnh bên $AD$ và $BC$ có độ dài bằng nhau}
{\True Hai góc kề một đáy của hình thang cân $ABCD$ bằng nhau}
{Hai đường chéo $AC$ và $BD$ vuông góc với nhau}
{\True Diện tích của hình thang cân $ABCD$ là $35\text{ cm}^2$}
\loigiai{
\textbf{a)} Đúng. Trong hình thang cân, hai cạnh bên luôn có độ dài bằng nhau.

\textbf{b)} Đúng. Đây là tính chất đặc trưng cơ bản của hình thang cân.

\textbf{c)} Sai. Trong hình thang cân, hai đường chéo bằng nhau chứ không nhất thiết vuông góc với nhau.

\textbf{d)} Đúng. Diện tích hình thang cân là: $S = \dfrac{(4 + 10) \cdot 5}{2} = 35\text{ (cm}^2\text{)}$.
}
\end{ex}


\setcounter{ex}{0}
\vspace{2mm}
\noindent\textbf{\large PHẦN III. CÂU HỎI TRẢ LỜI NGẮN}
\vspace{1mm}
\begin{ex}
Tìm số tự nhiên $x$ sao cho $24 \ \vdots \ x$ và $10 \le x < 15$. \shortans{12}
\loigiai{
Vì $24 \ \vdots \ x$ nên $x \in Ư(24) = \{1; 2; 3; 4; 6; 8; 12; 24\}$.\\Vì $10 \le x < 15$ nên $x = 12$.
}
\end{ex}

\begin{ex}
Một hình thang cân có độ dài hai đáy lần lượt là $6\text{ cm}$ và $10\text{ cm}$, chiều cao bằng $4\text{ cm}$. Tính diện tích của hình thang cân đó (theo đơn vị $\text{cm}^2$). \shortans{32}
\loigiai{
Diện tích hình thang cân là: $S = \dfrac{(6 + 10) \cdot 4}{2} = 32\text{ (cm}^2\text{)}$.
}
\end{ex}


\setcounter{ex}{0}
\vspace{2mm}
\noindent\textbf{\large PHẦN IV. TỰ LUẬN (BÀI TẬP VÀ BAREM ĐIỂM)}
\vspace{1mm}
\begin{ex} (1,0 điểm)
Thực hiện phép tính (Tính hợp lí nếu có thể):
\begin{enumerate}[label=\textbf{\alph*)}, leftmargin=0.8cm, itemsep=2pt]
\item (0,3 điểm) $6^2 : 4 + 2^3 \cdot 5$
\item (0,35 điểm) $47 \cdot 31 + 69 \cdot 47 - 200$
\item (0,35 điểm) $12 \cdot \{390 : [500 - (125 + 245)]\}$
\end{enumerate}
\loigiai{
\textbf{a)} \textbf{[0,3 điểm]}\\ $6^2 : 4 + 2^3 \cdot 5 = 36 : 4 + 8 \cdot 5 = 9 + 40 = 49$.

\textbf{b)} \textbf{[0,35 điểm]}\\ $47 \cdot (31 + 69) - 200 = 47 \cdot 100 - 200 = 4700 - 200 = 4500$.

\textbf{c)} \textbf{[0,35 điểm]}\\ $12 \cdot \{390 : [500 - 370]\} = 12 \cdot \{390 : 130\} = 12 \cdot 3 = 36$.


}
\end{ex}
\vspace{2mm}

\begin{ex} (1,0 điểm)
Tìm số tự nhiên $x$, biết:
\begin{enumerate}[label=\textbf{\alph*)}, leftmargin=0.8cm, itemsep=2pt]
\item (0,3 điểm) $x + 35 = 98$
\item (0,35 điểm) $3x - 45 = 60$
\item (0,35 điểm) $(6 - x)^3 + 11 = 19$
\end{enumerate}
\loigiai{
\textbf{a)} \textbf{[0,3 điểm]}\\ $x = 98 - 35 \Rightarrow x = 63$.

\textbf{b)} \textbf{[0,35 điểm]}\\ $3x = 60 + 45 \Rightarrow 3x = 105 \Rightarrow x = 105 : 3 = 35$.

\textbf{c)} \textbf{[0,35 điểm]}\\ $(6 - x)^3 = 19 - 11 \Rightarrow (6 - x)^3 = 8 \Rightarrow (6 - x)^3 = 2^3 \Rightarrow 6 - x = 2 \Rightarrow x = 4$.


}
\end{ex}
\vspace{2mm}

\begin{ex} (0,5 điểm)
Bác Minh muốn lát nền cho một căn phòng hình chữ nhật có chiều dài $16\text{ m}$, chiều rộng $5\text{ m}$. Loại gạch lát nền được sử dụng là gạch hình vuông có cạnh dài $40\text{ cm}$. Hỏi bác Minh phải sử dụng bao nhiêu viên gạch (coi mạch vữa không đáng kể)?
\loigiai{
Diện tích căn phòng là: $16 \times 5 = 80\text{ (m}^2\text{)} = 800\,000\text{ (cm}^2\text{)}$.\\Diện tích một viên gạch lát nền là: $40 \times 40 = 1600\text{ (cm}^2\text{)}$.\\Số viên gạch bác Minh cần sử dụng là: $800\,000 : 1600 = 500\text{ (viên)}$.
}
\end{ex}
\vspace{2mm}

\begin{ex} (0,5 điểm)
Chứng tỏ rằng biểu thức sau chia hết cho 21:\\$A = 1 + 4 + 4^2 + 4^3 + \dots + 4^{2024}$
\loigiai{
Tổng $A$ có $2024 - 0 + 1 = 2025$ số hạng. Vì $2025 \ \vdots \ 3$, ta nhóm 3 số hạng liên tiếp vào một nhóm:
\begin{align*}
A &= (1 + 4 + 4^2) + 4^3(1 + 4 + 4^2) + \dots + 4^{2022}(1 + 4 + 4^2)\\
&= (1 + 4 + 16) + 4^3 \cdot (1 + 4 + 16) + \dots + 4^{2022} \cdot (1 + 4 + 16)\\
&= 21 + 4^3 \cdot 21 + \dots + 4^{2022} \cdot 21\\
&= 21 \cdot (1 + 4^3 + \dots + 4^{2022})
\end{align*}
Vì tích có thừa số 21 nên $A \ \vdots \ 21$ (đpcm).
}
\end{ex}
\vspace{2mm}


\begin{center}
\rule{4cm}{0.4pt}\\
\textit{--- HẾT ---}
\end{center}
\end{document}
